\chapterbackmatter{Supplementary Exercises}

\begin{enumerate}
\SupplementaryExercise{1}{Connected Moments, M{\o}bius Inversion, and Independence}
  {(Inspired by D. Harlow, \emph{Quantum Field Theory I}, 2024, Section 7.5, Eqs. 7.47--7.49 | D. Harlow, \emph{Quantum Field Theory I}, 2024, Section 7.5 | D. Harlow, \emph{Quantum Field Theory I}, 2024, Sections 7.4--7.5.)}
  {harlow-qft1-linked-2024-curated-parent-1}
\begin{enumerate}
\item
For four not-necessarily-centered random variables \(X_1,\ldots,X_4\),
write the full fourth moment as a sum over set partitions into connected
moments.  Invert the relation explicitly to obtain
\(\langle X_1X_2X_3X_4\rangle_c\).
\item
If moments and connected moments obey
\[
 m(A)=\sum_{\pi\in\Pi(A)}\prod_{B\in\pi}\kappa(B),
\]
prove the inverse formula
\[
 \kappa(A)=\sum_{\pi\in\Pi(A)}
 (-1)^{|\pi|-1}(|\pi|-1)!\prod_{B\in\pi}m(B).
\]
\item
Let two collections of variables \(X\) and \(Y\) be statistically
independent.  Using a joint generating function, prove that every connected
moment containing at least one \(X\) and at least one \(Y\) vanishes.
Relate this result to cluster decomposition.
\end{enumerate}

\SupplementaryExercise{2}{Generating Connected Contributions and Their Volume Scaling}
  {(Inspired by D. Harlow, \emph{Quantum Field Theory I}, 2024, Section 7.4 | D. Harlow, \emph{Quantum Field Theory I}, 2024, Section 7.5 | D. Harlow, \emph{Quantum Field Theory I}, 2024, Sections 7.4--7.5.)}
  {harlow-qft1-linked-2024-curated-parent-2}
\begin{enumerate}
\item
Suppose each connected diagram type \(C\) has value \(V_C\), and a
disconnected diagram contains \(m_C\) identical copies of each type.
Show that its combinatorial weight is
\(\prod_C V_C^{m_C}/m_C!\), and prove that the sum over all diagrams is
\(\exp(\sum_CV_C)\).
\item
Let \(Z[J]\) be normalized by \(Z[0]=1\), and define \(W[J]=\log Z[J]\).
Show through fourth order in \(J\) that functional derivatives of \(W\)
at \(J=0\) give connected correlation functions, while derivatives of
\(Z\) give the full functions.
\item
In a translationally invariant theory in a periodic box of volume \(V\),
show that a connected momentum-space amplitude carries one momentum
Kronecker delta and scales as \(V\), whereas a contribution with \(C\)
disconnected components scales as \(V^C\).  Explain why the logarithm is
extensive.
\end{enumerate}

\SupplementaryExercise{3}{Vacuum Cancellation, Connectedness, and Momentum Delta Functions}
  {(Inspired by D. Tong, \emph{Lectures on Quantum Field Theory}, 2006, Section 3.7.1 | D. Tong, \emph{Lectures on Quantum Field Theory}, 2006, Sections 3.5.4 and 3.7.1 | D. Tong, \emph{Lectures on Quantum Field Theory}, 2006, Sections 3.4.1 and 3.5.4.)}
  {tong-qft-connected-2006-curated-parent-1}
\begin{enumerate}
\item
Use the Gell-Mann--Low ratio for an \(n\)-point function to prove that all
vacuum bubbles cancel between numerator and denominator.  Distinguish a
vacuum bubble from a diagram disconnected only with respect to external
points.
\item
Give graph-theoretic definitions of a connected graph and a
one-particle-irreducible graph.  Construct a connected two-point graph
that becomes disconnected when one internal line is cut, and explain why
it contributes to a connected correlator.
\item
For a diagram with \(V\) vertices, \(I\) internal lines, \(L\) loop
momenta, and \(C\) connected components, prove
\(V-I+L=C\).  Use momentum conservation at every vertex to show that
exactly \(C\) independent momentum delta functions remain.
\end{enumerate}

\SupplementaryExercise{4}{Locality, Singular Support, and Cluster Factorization}
  {(Inspired by D. Tong, \emph{Lectures on Quantum Field Theory}, 2006, Sections 1.1.4 and 3.4.1 | D. Tong, \emph{Lectures on Quantum Field Theory}, 2006, Section 1.1.4 | D. Tong, \emph{Lectures on Quantum Field Theory}, 2006, Section 3.5.4.)}
  {tong-qft-connected-2006-curated-parent-2}
\begin{enumerate}
\item
Fourier transform a local interaction
\[
 H_I(t)=\frac{\lambda}{4!}\int d^3x\,\phi(\mathbf x,t)^4.
\]
Show that each term in momentum space contains one and only one
three-momentum-conservation delta function.
\item
Consider
\[
 H_{\rm biloc}=\frac12\int d^3x\,d^3y\,
 K(\mathbf x-\mathbf y)\phi(\mathbf x)^2\phi(\mathbf y)^2.
\]
Find its momentum-space kernel.  Determine when it contains more singular
support than one total-momentum delta function, and analyze the special
case \(K(\mathbf r)=K_0\).
\item
For distinguishable particles, expand the \(4\to4\) \(S\)-matrix into
connected pieces.  Then separate particles \(1,2\) far from \(3,4\) and
show directly that the surviving terms combine into
\(S_{1'2',12}S_{3'4',34}\).
\end{enumerate}

\SupplementaryExercise{5}{Wick's Theorem for Three Fields}
  {(Adapted from B. Mistlberger and K. Zhou, Stanford PHYS-330, Problem Set 4, 2022, Problem 1.)}
  {zhou-stanford-phys330-ps4-2022-curated-parent-1}
\begin{enumerate}
\item For a free scalar, compare time ordering, normal ordering, and the
Feynman two-point function.  Prove
\(T\{\phi(x)\phi(y)\}=:\!\phi(x)\phi(y)\!:+D_F(x-y)\).
\item Then prove Wick's theorem for
\(T\{\phi(x_1)\phi(x_2)\phi(x_3)\}\) by separating every field into
creation and annihilation parts, displaying all single contractions.
\end{enumerate}

\SupplementaryExercise{6}{Cubic-Theory Vacuum Graphs and Normalized Correlators}
  {(Adapted from B. Mistlberger and K. Zhou, Stanford PHYS-330, Problem Set 4, 2022, Problem 2.)}
  {zhou-stanford-phys330-ps4-2022-curated-parent-2}
\begin{enumerate}
\item
For
\(\mathcal L_{\rm int}=+\lambda\phi^3/3!\), evaluate
\(\langle0|T\exp[-i\int dt\,H_I(t)]|0\rangle\) through order
\(\lambda^2\).  Identify the connected vacuum terms and their symmetry
factors.
\item
Draw and count every vacuum topology at orders \(\lambda^3\) and
\(\lambda^4\), assigning a symmetry factor to each.  Explain from
pairing parity why the third-order list is empty.
\item
Draw every two-point diagram through order \(\lambda^4\), including
disconnected vacuum factors, and assign its Wick-contraction symmetry
factor.  Then expand the normalized numerator and denominator and show
that all vacuum components cancel.
\item
Explain the difference between (i) every interaction vertex being
connected to at least one external point and (ii) every external point
belonging to one connected component.  Give a four-point example that
satisfies (i) but not (ii), and show how the connected four-point
function removes it.
\item
In cubic theory at order \(\lambda^2\), compute the symmetry factors of
the connected theta vacuum graph (three lines joining the two vertices)
and of the graph with one line joining the vertices and a tadpole at each
vertex.  Derive the answer by counting contractions.
\end{enumerate}

\SupplementaryExercise{7}{Coherent States, Completeness, and In--Out Vacua}
  {(Adapted from A. Guth, MIT 8.323, Problem Set 5, 2008, Problem 3.)}
  {mit-823-s08-ps5-curated-parent-1}
\begin{enumerate}
\item
For one oscillator, construct the normalized eigenstate
\(a\ket z=z\ket z\), expand it in number states, and compute
\(\braket{z_2}{z_1}\).
\item
Prove
\[
 \int_{\mathbb C}\frac{d^2z}{\pi}\ket z\bra z=\mathbf1
\]
for normalized oscillator coherent states.  Explain why an overcomplete
family can resolve the identity even though its members are not
orthogonal.
\item
With \(q=(a+a^\dagger)/\sqrt2\) and
\(p=(a-a^\dagger)/(i\sqrt2)\), compute
\(\langle q\rangle,\langle p\rangle,\Delta q,\Delta p\) in \(\ket z\)
and verify minimum uncertainty.
\item
For a normalized field coherent state \(\ket\lambda\), find the generating
function \(\langle e^{tN}\rangle\) of the total particle number
\(N=\int dq\,a_q^\dagger a_q\).  Derive its mean, variance, and
probability distribution.
\item
Define \(D(z)=\exp(za^\dagger-z^*a)\).  Use the
Baker--Campbell--Hausdorff formula to normal order \(D(z)\), prove it is
unitary, and show \(D(z)\ket0=\ket z\).
\item
A classical source produces the relation
\[
 a_{\rm out}(\mathbf p)=a_{\rm in}(\mathbf p)+\alpha(\mathbf p).
\]
Express the in-vacuum as a coherent state built on the out-vacuum and
compute the no-particle probability.
\end{enumerate}

\SupplementaryExercise{8}{Classical Radiation from a Coherent Impulse}
  {(Adapted from B. Mistlberger and K. Zhou, Stanford PHYS-330, Problem Set 3, 2022, Problem 2.)}
  {zhou-stanford-phys330-ps3-2022-curated-parent-1}
\begin{enumerate}
\item
A free real scalar field couples to
\(J(\mathbf x,t)=\delta(t)j(\mathbf x)\).  Integrate the equations of
motion across the impulse, determine the shift of \(a(\mathbf p)\), and
show that the post-impulse state is coherent.
\item
For a free scalar field in a multimode coherent state, compute
\(\langle\phi(x)\rangle\) and prove that it obeys the classical
Klein--Gordon equation.  Show that the connected two-point function is
the same as in the vacuum.
\item
For the impulsive source of S.4.8(a), compute the mean occupation density,
the total mean number, and the probability of producing no quanta.
State the condition for the coherent state to be normalizable.
\end{enumerate}

\SupplementaryExercise{9}{Relativistic Propagators and Causal Support}
  {(Adapted from B. Mistlberger and K. Zhou, Stanford PHYS-330, Problem Set 3, 2022, Problem 1.)}
  {zhou-stanford-phys330-ps3-2022-curated-parent-2}
\begin{enumerate}
\item Starting from the free-field expansion, obtain
\[
 D(z)=\int\frac{d^4p}{(2\pi)^3}\,
 \delta(p^2+m^2)\theta(p^0)e^{ip\cdot z}.
\]
\item Prove that \(D(z)\) is Lorentz invariant.
\item Evaluate it at spacelike separation in terms of a modified Bessel
function and extract the large-distance exponential behavior.
\item Use the result to prove that the scalar commutator vanishes outside
the light cone.
\item Couple the Klein--Gordon field to a c-number source and derive its
inhomogeneous equation of motion.
\item Solve that equation using an arbitrary Green function plus a
homogeneous solution.
\item Define the retarded, advanced, and Feynman Green functions by their
momentum-contour prescriptions and write their position-space support
properties.
\item Show explicitly that the retarded solution vanishes at spacelike
separation from the source.
\item Show that the Feynman function need not vanish there, and explain
why this correlation does not permit acausal signaling.
\end{enumerate}

\SupplementaryExercise{10}{Separated Sources and Additivity of \(W\)}
  {(Inspired by B. Mistlberger and K. Zhou, Stanford PHYS-330, Problem Set 3, 2022, Problems 1--2.)}
  {zhou-stanford-phys330-ps3-2022-curated-parent-3}
Let \(J=J_A+J_B\), with the two sources supported in widely separated
regions and with negligible connected propagation between them.  Show
that \(Z[J]\to Z[J_A]Z[J_B]\) and hence
\(W[J]\to W[J_A]+W[J_B]\).  Relate this to independent experiments.
\end{enumerate}
